% Folder 79, batch 11 — archivists' pages 201–220 (author's pp. 104–113).
% Pass: Opus 5.5 (claude-opus-5-5), 2026-10-03 — first pass, unchecked against the pages by a human.
%
% Notation fixed for this batch: his dual set, written with a superscript
% "v" (sometimes looking like a star), is S^{\vee}; where he writes the
% accent over the letter it is \check{S}. His script R is \mathcal{R}, the
% symmetric group \mathfrak{S}_S, his script G is \mathcal{G}, his script H
% \mathcal{H}. Star sets Et_{A_0}(s) are \mathrm{Et}.
\documentclass[11pt,a4paper]{article}
\input{../preamble/grothendieck}

\folder{79}
\batch{11}
\pages{201}{220}
\foldertitle{2-polyèdres réguliers triangulés, et modules libres de rangs 2 sur les anneaux $\Lambda=\mathbb{Z}/n\mathbb{Z}$ et $\Lambda=\mathbb{Z}$ : notes manuscrites (s.d.).}
\dating{[à partir de 1980]}
\watermark{Édition de démonstration}

\begin{document}

\page{201}
\note{La page reprend une discussion commencée avant ce lot : « les descriptions précédentes » et les formules (21) et (22) renvoient aux pages qui précèdent.}
Les descriptions précédentes de la notion de graphe \uncertain{ont} l'avantage
principal d'être « \uncertain{intrinsèques} », \emph{et} de séparer deux
composantes de la structure :

\[
(22)\qquad
\begin{cases}
\text{a) la donnée } (S, S^{\vee}, \mathcal{R}_{A} \subset S \times S^{\vee}) \\
\text{b) la donnée } \varphi_0 : S \xrightarrow{\ \sim\ } S^{\vee} .
\end{cases}
\]

Rien n'empêche d'ailleurs d'étudier de telles structures, sans supposer que la
relation a) soit « bisignante » — i.e. en ne \uncertain{faisant} sur la
\uncertain{correspondance} aucune hypothèse particulière. L'\add{auto}dualité
(21) reste valable.

Le jeu nouveau, c'est qu'on peut partir d'une structure a)
\[
(23)\qquad (S, S^{\vee}, \mathcal{R}_{A} \subset S \times S^{\vee}),
\]
sans donnée d'un $\varphi_0$, — donc une structure en un sens plus faible
qu'une structure de graphe, — pouvant s'appeler structure de « prégraphe » —
structure considérée comme « préliminaire » à celle de graphe (bisignant
strict). Il se pose alors la question d'étudier toutes les structures de
graphe correspondantes, i.e. \emph{toutes} les applications bijectives
\struck{(sur $\mathcal{R}_0$}
\[
(24)\qquad \varphi : S \xrightarrow{\ \sim\ } S^{\vee}
\]
telles que $\varphi$ soit « $\mathcal{R}_A$-symétrique et antiréflexive »
(on dit que $\varphi$ est une $\mathcal{R}_A$-jauge) ; \struck{i.e. les}
« $\mathcal{R}_A$-jauges ». On s'intéresse surtout au cas, bien sûr, où cet
ensemble n'est pas vide — et le $\varphi_0$ \uncertain{du travail} joue le
rôle d'une « origine » dans l'ensemble

\marginal{antiréflexivité : \ill{} \ill{} \ill{} \ill{} \ill{} des bandes, \ill{} \ill{}}

\page{202}
\note{Pagination de l'auteur : p. 104.}
\[
(25)\qquad \Pi = \Pi_{\mathcal{R}} = \{\varphi : S \simeq S^{\vee} \mid \varphi \text{ est } \mathcal{R}\text{-symétrique}\}
\]
des $\mathcal{R}$-jauges. \struck{On a donc \ill{}} On voit que (pour un tel
choix $\varphi_0$
\[
(26)\qquad \varphi_0 : S \simeq S^{\vee}, \qquad \varphi_0 \in \Pi\ ),
\]
les applications bijectives $\varphi : S \simeq S^{\vee}$ correspondent
biunivoquement aux permutations $\lambda \in \mathfrak{S}_S$, via
\[
(27)\qquad \lambda \longmapsto \varphi = \varphi_{\lambda} = \varphi_0 \circ \lambda, \qquad \lambda \in \mathfrak{S}_S,
\]
ou encore aux $\lambda^{\vee} \in \mathfrak{S}_{S}$\note{l'indice est écrit
$S$ ; le sens demande $S^{\vee}$, comme dans la ligne suivante.}
\[
(27^{\vee})\qquad \lambda^{\vee} \longmapsto \varphi = \varphi_{\lambda^{\vee}} = \text{\struck{$\varphi$}}\ \lambda^{\vee} \circ \varphi_0, \qquad \lambda^{\vee} \in \mathfrak{S}_{S^{\vee}},
\]
la relation entre $\lambda$ et $\lambda^{\vee}$ étant donnée par
\[
(28)\qquad
\begin{cases}
\lambda^{\vee} = \varphi_0(\lambda) = \varphi_0 \lambda \varphi_0^{-1} \\
\lambda = \varphi_0^{-1}(\lambda^{\vee}) = \varphi_0^{-1} \lambda^{\vee} \varphi_0 .
\end{cases}
\]
Ceci dit, si $\varphi$ correspond à $\lambda$, pour que $\varphi$ soit une
$\mathcal{R}$-jauge, il faut et il suffit que l'on ait
\[
(29)\qquad
\begin{cases}
\text{la relation } (s, \underbrace{\varphi_0(\lambda(t))}_{\lambda^{\vee}\varphi_0(t)}) \in \mathcal{R} \text{ est symétrique en } s,t, \\
\text{i.e. la relation } \{s, \lambda(t)\} \in A_0 \text{ symétrique en } (s,t) \\
\text{et elle est antiréflexive, i.e. } \forall s \in S, \text{ on a } \{s, \lambda(s)\} \notin A_0 .
\end{cases}
\]
La symétrie s'écrit aussi sous la forme
\[
(29')\qquad \forall (s,t) \in S \times S, \text{ on a } \{s,t\} \in A_0 \Longleftrightarrow \{\lambda^{-1}s, \lambda t\} \in A_0
\]
(il suffit que $\forall s,t \in S \times S$,
$\{s,t\} \in A_0 \Longrightarrow \{\lambda^{-1}s, \lambda t\} \in A_0$)
\struck{ou encore}
— ou encore
\[
(29'')\qquad \forall s \in S, \text{ on a } \lambda(\underbrace{\mathrm{Et}_{A_0}(s)}_{\varphi_0(s)}) = \underbrace{\mathrm{Et}(\lambda^{-1}s)}_{\varphi_0(\lambda^{-1}s)}
\]
— sous cette forme, est bien mise en évidence l'action

\marginal{antiréflexivité \ill{} \ill{}, \ill{} exceptées \ill{} graphes \ill{}}

\page{203}
de $\lambda$ sur les étoiles\add{$E = \mathcal{R}_{s^{\vee}}$}
\struck{(elle se fait comme de juste par transport des structures}

\struck{(30)\quad $\lambda(E) = \lambda({}^{s}\mathcal{R}_{s^{\vee}}) =$} ).

Notons que le couple
\[
(30)\qquad (\lambda, \lambda^{\vee}) \in \mathfrak{S}_S \times \mathfrak{S}_{S^{\vee}}
\]
défini par \uncertain{deux} couples $(\varphi_0, \varphi)$ de
\struck{\ill{}} $\mathcal{R}$-jauges, satisfait la condition
\[
(31)\qquad \text{\struck{$(s, s^{\vee}) \in \mathcal{R} \Longleftrightarrow \ldots$}}\quad
(s, \check{s}) \in \mathcal{R} \Longleftrightarrow (\lambda^{-1}s, \lambda^{\vee} s^{\vee}) \in \mathcal{R}
\]
qui \struck{en fait} (pour $\varphi_0$ fixé) en fait \uncertain{(29')} sur les
$\varphi = \lambda^{\vee}\varphi_0 = \varphi_0\lambda : S \to S^{\vee}$
\ill{}, qui exprime la condition définie par \struck{\ill{}} $\lambda$.
\note{Ce passage, raturé et surchargé (plusieurs formes successives de (31)
biffées), n'est lu que partiellement.}

\marginal{Digression}
\emph{Exemples.} Soient $\Lambda$ un anneau, $(M, M^{\vee})$\add{\uncertain{noté $M^{\vee}$}}
un couple de $\Lambda$-modules libres de rang 2, accouplés par une forme
bilinéaire identifiant chacun d'eux au dual de l'autre. Considérons les
\struck{relat} ensembles
\[
(32)\qquad
\begin{cases}
M^{*} \subset M \\
M^{\vee *} \subset M^{\vee} \\
S = M^{*}/\pm 1, \quad S^{\vee} = M^{\vee *}/\pm 1
\end{cases}
\]
(ens. des $x \in M$ (resp. $\check{x} \in M^{\vee}$) faisant partie d'une base),
et la relation
\[
(33)\qquad \mathcal{R} \subset S \times \check{S},\quad
\mathcal{R} = \{(s, \check{s}) \mid \exists (x, \check{x}) \in M^{*} \times M^{\vee *} \text{ t.q. } s = \dot{x},\ \check{s} = \dot{\check{x}},\ \langle x, \check{x} \rangle = 1\}.
\]

\note{La suite de la page est barrée de longs traits obliques ; elle se lit
néanmoins.}
\struck{Nous nous restreindrons aux applications}
\struck{$\varphi_0 : S \simeq S^{\vee}$ induites par des applications linéaires
(notées par la même lettre)}
\[
\text{\struck{$(34)\quad \varphi_0 : M \xrightarrow{\ \sim\ } M^{\vee}$ \quad $\Lambda$-linéaires,}}
\]
\struck{qui correspondent donc aux formes bilinéaires non dégénérées
$M \times M \to \Lambda$, $\psi_0(x,y) = (x, \varphi_0' y)$),}
\struck{satisfaisant la condition}
\struck{qui sont \emph{alternées}}
\[
\text{\struck{$\psi_0(x,x) = 0$ \quad i.e. \quad $(x, \varphi_0(x)) = 0$,}}
\]
\struck{et qui sont donc en corr. 1-1 avec les $e^{\vee} \in \Lambda^2 M^{\vee} = \det(M)^{\vee}$,}

\page{204}
\note{Pagination de l'auteur : p. 105.}
On l'obtient en relation avec la structure de graphe\add{$(S, A_e)$} sur
$S$, associée au choix d'une base\add{de $\det M = \Lambda^2 M$} — ou
« jauge » de $M$ —
\[
(34)\qquad e \in \det(M)^{*}
\]
par la \struck{condition} formule
\[
(35)\qquad A_e = \bigl\{\{s,t\} \in \mathfrak{P}_2(S) \bigm| \exists\, x, y \in M^{*} \text{ t.q. } \dot{x} = s,\ \dot{y} = t,\ x \wedge y = e\ \text{\struck{\ill{}}}\bigr\}
\]
(on \uncertain{dispose} que de la structure \uncertain{symplectique} de $M$
$\alpha = \{e, -e\}$ définie par $e$). Posant
\[
(36)\qquad
\begin{cases}
x \wedge y = \psi_e(x,y)\, e \\
\psi_e \text{ forme bilinéaire alternée non dég. définie par } e ,
\end{cases}
\]
on voit donc que l'on a, pour $s \in S$, $s = \dot{x}$,
\[
(37)\qquad \mathrm{Et}_{A_e}(s) = \{\dot{y} \mid y \in M^{*},\ \psi_e(x,y) = \pm 1\} .
\]
Considérons alors l'application linéaire
\[
(38)\qquad
\begin{cases}
\Phi_e : M \xrightarrow{\ \sim\ } \check{M} \quad \text{caractérisée par} \\
\langle x, \Phi_e(y) \rangle = \psi_e(x,y) \quad \text{i.e.} \quad x \wedge y = \langle x, \Phi_e(y) \rangle \cdot e .
\end{cases}
\]
Elle induit, par passage au quotient,
\struck{(39) donc}
$\varphi_e : S \to \check{S}$, et on a, pour $s \in S$,
\[
(40)\qquad \mathrm{Et}_{A_e}(s) = \{t \in S \mid (s, \varphi_e(t)) \in \mathcal{R}_A\} ,
\]
où
\[
(41)\qquad
\begin{aligned}
\mathcal{R}_A \subset S \times S^{\vee}
&= \{(\dot{x}, \dot{x}^{\vee}) \in S \times S^{\vee} \mid \langle x, x^{\vee} \rangle = \pm 1\} \\
&= \{(s, \check{s}) \in S \times \check{S} \mid \langle s, \check{s} \rangle = \pm 1\}
\end{aligned}
\]
\marginal{(42) où $\langle\ ,\ \rangle : S \times \check{S} \to \Lambda/\pm\mathrm{id}_{\Lambda}$
\uncertain{est induit} par passage au quotient \uncertain{par}
$M^{*} \times M^{\vee *} \to \Lambda$}
\struck{$S$-graphes, on voit donc que} les structures de graphe $(S, A_e)$
strictes, i.e. qu'on a
\[
(43)\qquad
\begin{cases}
\forall (x, x') \in M^{*} \times M^{*}, \text{ la relation} \\
\quad [\forall y \in M^{*}\quad x \wedge y \in (\det M)^{*} \Longleftrightarrow x' \wedge y \in (\det M)^{*}] \\
\qquad\qquad \Longrightarrow x' \wedge y \in \{\pm x \wedge y\} \\
\text{implique que } x' \in \pm x,\ \text{\struck{\ill{}}}
\end{cases}
\]
ce qui, pour $\Lambda$ ½ local\note{« ½ local » : lecture du signe
incertaine ; peut-être « semi-local ».}, signifie simplement qu'en

\page{205}
n'est pas dans le cas octaédral généralisé
\[
(44)\qquad 2\Lambda = \{0, 2\}, \quad 2 \cdot 1_{\Lambda} \neq 0
\]
et dans le cas $\Lambda = \mathbb{Z}/n\mathbb{Z}$ signifie simplement que
$n \neq 4$), alors il y a corr. 1-1 entre étoiles de $(S, A_e)$, et éléments
de $S^{\vee}$, via la relation $\mathcal{R}$ de (41), et la structure de graphe
$(S, A_e)$ est associée à la structure
\[
(45)\qquad (S, \check{S}, \mathcal{R} \subset S \times \check{S},\ \varphi_e : S \xrightarrow{\ \sim\ } \check{S}) .
\]
Notons maintenant que le groupe
\[
(46)\qquad Z'_{\Lambda} = \Lambda^{*}/\pm 1
\]
opère sur $S$ et sur $S^{\vee}$ par symétrie, \struck{et} par des
multiplications\add{provenant de $\lambda \in \Lambda^{*}$} $(\dot{\lambda} \in Z'_{\Lambda})$,
opérations que je vais noter \struck{\ill{}}
\[
(47)\qquad \dot{\lambda} \cdot \dot{x} \overset{\text{déf}}{=} (\lambda x)^{\cdot}, \qquad
\dot{\lambda} \cdot (x^{\vee})^{\cdot} \overset{\text{déf}}{=} (\lambda \check{x})^{\cdot} .
\]
\struck{Ce sont} Ces opérations sont des opérations \emph{libres}, et on a
\[
(48)\qquad
\boxed{\begin{array}{l}
\forall (s, \check{s}) \in S \times S^{\vee}, \text{ et } \dot{\lambda} \in Z'_{\Lambda}, \text{ on a} \\
(s, s^{\vee}) \in \mathcal{R} \Longleftrightarrow (\dot{\lambda} s, \dot{\lambda}^{-1} \check{s}) \in \mathcal{R}
\end{array}}
\]
Notons aussi que
\[
(49)\qquad \varphi_e = \varphi_{-e} : S \simeq \check{S},
\]
\struck{donc pour $\lambda \in \Lambda^{*}$, $\varphi_{\lambda e}$ ne dépend que de}
i.e. $\varphi_e$ ne dépend que de l'élément
\[
(50)\qquad \alpha = \{e, -e\} \in \pi \overset{\text{déf}}{=} (\det M)^{*}/\pm 1 \underset{e \mapsto \check{e}^{-1}}{\simeq} (\det M^{\vee})^{*}/\pm 1
\]
défini par $e$, de sorte qu'on a en fait
\[
(51)\qquad (\varphi_{\alpha})_{\alpha \in \pi}, \quad \text{famille de bijections indexée par } \pi .
\]
\marginal{cf. (54) provision}
Notons que \struck{$\pi$ est un torseur}
\[
(52)\qquad \pi \text{ est un torseur sous } Z'_{\Lambda},
\]
par des opérations notées $\alpha \mapsto T_{\dot{\lambda}}\alpha$
($\dot{\lambda} \in Z'_{\Lambda}$), déduites, par passage au quotient, des
opérations $e \mapsto \lambda . e$.
\marginal{cf. (55) ; (52) provision}

\page{206}
\note{Pagination de l'auteur : p. 106.}
\struck{On considère alors $\varphi_{\dot{\lambda} e}$}

Notons que
$\varphi_{\dot{\lambda} e}(\dot{x}) \in \check{S} = M^{\vee *}/\pm 1$ est défini
comme $\Phi_{\lambda e}(x)^{\cdot}$, où $\Phi_{\lambda e}(x) \in M^{\vee}$ défini
par (38),
\struck{$\varphi_{\lambda e}(x)$ défini par}
\[
\langle y, \Phi_{\lambda e}(x) \rangle . (\lambda e) = y \wedge x .
\]
Puis comparons avec
\[
\langle y, \Phi_e(x) \rangle\, e = y \wedge x ,
\]
on trouve
\[
\Phi_{\lambda e}(x) = \lambda^{-1} \Phi_e(x) = \Phi_e(\lambda^{-1} x)
\]
sur les éléments correspondants de $S^{\vee}$
soit, en passant aux éléments correspondants de $S^{\vee}$,
\[
(53)\qquad \varphi_{\dot{\lambda} . \alpha}(s)\ \text{\struck{\ill{}}} = \dot{\lambda}^{-1} \varphi_e(s) = \varphi_e(\dot{\lambda}^{-1} s)
\qquad \forall s \in S^{\vee},\ \dot{\lambda} \in Z'_{\Lambda} .
\]
\marginal{(53) provision, cf. (55)}
Si on avait utilisé $\check{\pi} = (\det M^{\vee})^{*}$ pour indexer les
$\varphi : S \to S^{\vee}$ que nous considérons, ce qui était plutôt plus
naturel (puisque les applications linéaires
\struck{$M \to \check{M}$} $\Phi : M \to M^{\vee}$ définies en sens
\uncertain{précis}, sont \uncertain{les} « mêmes » que les applications
\uncertain{constantes}\note{lecture douteuse ; le sens attendu serait
« bilinéaires ».} bilinéaires (alternées) $M \times M \to \Lambda$, i.e. aux
formes linéaires $\Lambda^2 M \to \Lambda$, et s'identifient donc aux
\ill{} : aux éléments $\check{\alpha}$ $(= \alpha^{-1})$ de
\[
(\det M^{\vee})' \simeq (\det M)^{\vee}
\]
Ainsi on va considérer plutôt qu'en défaut
\[
(54)\qquad \pi = \pi_M = (\det M^{\vee})^{*}/\pm \mathrm{id}_{M^{\vee}},
\]
sur lequel on fait opérer $Z'_{\Lambda} = \Lambda^{*}/\pm 1$ par l'opération
habituelle — qui correspond \struck{\ill{}} : la multiplication d'une forme
bilinéaire — et (53), \struck{\ill{}} remplacée par
\[
(55)\qquad
\boxed{\varphi_{\dot{\lambda} . \alpha}(s) = \dot{\lambda}\, \varphi_{\alpha}(s) = \varphi_{\alpha}(\dot{\lambda} s)}
\qquad \forall \alpha \in \pi = (\det M^{\vee})^{*}/\pm 1,\ \dot{\lambda} \in Z'_{\Lambda} = \Lambda^{*}/\{\pm 1\}
\]

\page{207}
\struck{4.} Appelons « prégraphe » un système
\[
(56)\qquad \Gamma = (S, S^{\vee}, \mathcal{R} \subset S \times S^{\vee})
\]
de deux ensembles $S, S^{\vee}$, et d'une relation entre eux — il sera dit
\emph{séparé} si les applications
\[
S \to \mathfrak{P}(S^{\vee}), \qquad S^{\vee} \to \mathfrak{P}(S)
\]
correspondantes sont injectives. Le symétrique de (56) est le prégraphe
\[
(57)\qquad \check{\Gamma} = (\check{S}, S, {}^{s}\mathcal{R} \subset S^{\vee} \times S) .
\]
On appelle « \emph{homothétie} » du prégraphe toute transformation faite de
paires de permutations
\[
(58)\qquad \lambda = (\lambda_S, \lambda_{S^{\vee}}) \in \mathfrak{S}_S \times \mathfrak{S}_{S^{\vee}}
\]
satisfaisant la condition
\[
(59)\qquad
\begin{cases}
(\lambda_S, \lambda_{S^{\vee}}^{-1})(\mathcal{R}) = \mathcal{R} \quad \text{i.e.} \\
\forall (s, s^{\vee}) \in S \times S^{\vee}, \text{ on a } ((s, s^{\vee}) \in \mathcal{R} \Longleftrightarrow (\lambda s, \lambda^{-1} s^{\vee}) \in \mathcal{R}),
\end{cases}
\]
\marginal{cf. (54) provision}
\note{Dans la marge de la p.~205, à côté de (50), l'auteur renvoie déjà
« cf. (59) provision ».}
\uncertain{elles} sont donc les éléments \struck{$(\lambda, \lambda')$}
$(u, u') \in \mathfrak{S}_S \times \mathfrak{S}_{S^{\vee}}$ tels que l'on ait
\[
(60)\qquad
\begin{cases}
(u, u'^{-1}) \in \mathrm{Aut}(S, S', \mathcal{R}) \quad \text{i.e.} \\
(u^{-1}, u') \in \mathrm{Aut}(S, S', \mathcal{R}) .
\end{cases}
\]
Notons que si $\mathcal{R}$ est séparée en $\check{S}$, alors
$\lambda_{S^{\vee}}$ est connu quand on connaît $\lambda_S$, et est défini à
partir de $\lambda_S$ par
\[
(61)\qquad
\begin{cases}
\text{\struck{$\lambda_{S^{\vee}}(s^{\vee})$}}\ {}^{s}\mathcal{R}(\lambda_{S^{\vee}}(s^{\vee})) = \lambda_S^{-1}(\mathcal{R}(s^{\vee})) \quad \text{ou} \\
{}^{s}\mathcal{R}(\lambda_{S^{\vee}}^{-1}(s^{\vee})) = \lambda_S(\mathcal{R}(s^{\vee}))
\end{cases}
\]
et inversement, si $\mathcal{R}$ est séparée en $S$, $\lambda_{S^{\vee}}$
détermine $\lambda_S$ —

\page{208}
\note{Pagination de l'auteur : p. 107.}
donc dans le cas bisignant, $\lambda_S$ et $\lambda_{S^{\vee}}$ se
déterminent mutuellement. Dans ce cas, on trouve des bijections canoniques
\[
(62)\qquad
\begin{array}{l}
\text{Homothéties de } \Gamma \text{ i.e.\ couples } \lambda_S, \lambda_{S^{\vee}} \text{ satisfaisant (59)} \\
\quad \simeq\ \underbrace{\text{automorphismes de la structure } {}^{s}\Gamma = (S, \check{S}, \mathcal{R})}_{\text{soit } Z} \\[4pt]
\quad \underset{\text{soit } Z_S}{\simeq} \{\lambda \in \mathfrak{S}_S \mid \underbrace{\lambda(\Sigma)}_{\text{transport de str.}} = \Sigma\} \\
\qquad \text{où } \Sigma \subset \mathfrak{P}(S) \text{ est l'image de } S^{\vee} \xrightarrow{{}^{s}\mathcal{R}} \mathfrak{P}(S) \\[4pt]
\quad \underset{\text{soit } Z_{S^{\vee}}}{\simeq} \{\lambda^{\vee} \in \mathfrak{S}_{S^{\vee}} \mid \underbrace{\lambda^{\vee}(\Sigma^{\vee})}_{\text{transp. de str.}} = \Sigma^{\vee}\} \\
\qquad \text{où } \Sigma^{\vee} \subset \mathfrak{P}(S^{\vee}) \text{ est l'image de } S \xrightarrow{\mathcal{R}} \mathfrak{P}(S^{\vee})
\end{array}
\]
\note{L'auteur écrit « $\Sigma^{\vee} \subset \mathfrak{P}(S)$ » ; le sens
demande $\mathfrak{P}(S^{\vee})$, comme dans la flèche qu'il écrit dessous.}
Sauf le premier des trois\add{4} ensembles, des autres ensembles sont munis
des structures de groupes évidentes, qui sont isomorphes par les applications
évidentes (63)

\begin{tikzcd}[column sep=small]
 & Z \arrow[dl, "{p_S \,\sim}"'] \arrow[dr, "p_{S^{\vee}}"] & \\
Z_S \arrow[d, hook] & & Z_{S^{\vee}} \arrow[d, hook] \\
\mathfrak{S}_S & & \mathfrak{S}_{S^{\vee}}
\end{tikzcd}

mais l'ensemble des homothéties apparaît comme le sous-ensemble un peu
vicieux de $\mathfrak{S}_S \times \mathfrak{S}_{S^{\vee}}$
\[
(64)\qquad \text{\struck{\ill{}}}\ \mathcal{H} = \{(p_S(\lambda), p_{S^{\vee}}(\lambda^{-1})) \in \mathfrak{S}_S \times \mathfrak{S}_{S^{\vee}} \mid \lambda \in Z\},
\]
qui fait que ce n'\emph{est pas} a priori un sous-groupe.

\page{209}
Si p.\ ex.\ $\mathcal{R}$ est le graphe d'une \emph{bijection} entre $S$ et
$S^{\vee}$, alors on a $Z_S = \mathfrak{S}_S$, $Z_{S^{\vee}} = \mathfrak{S}_{S^{\vee}}$,
et identifiant $S^{\vee}$ à $S$ par $\mathcal{R}$, d'où
$\mathcal{H} \subset \mathfrak{S}_S \times \mathfrak{S}_S$,
$\mathcal{H}$ s'identifie à l'antidiagonale
\[
\mathcal{H} = \{(u, u^{-1}) \in \mathfrak{S}_S \times \mathfrak{S}_S \mid u \in \mathfrak{S}_S\},
\]
qui n'est pas un sous-groupe \struck{si $\operatorname{card} S \geq 3$}
\uncertain{dès que} $\mathfrak{S}_S$ n'est pas commutatif, i.e.
$\operatorname{card} S \geq 3$. La définition\add{des homothéties} par
paires commutantes, pour \uncertain{être} des conditions
\uncertain{analogues} \ill{} précédentes.

C'est par simplement \uncertain{par} un souci de notation — ainsi, rien ne
nous empêcherait de modifier la définition en \uncertain{appelant}
commutantes des couples $(\lambda, \lambda^{\vee})$ qualifiés
d'homothéties, en les remplaçant par les couples $(\lambda^{-1}, \lambda^{\vee})$ — ou
$(\lambda, \lambda^{\vee -1})$ — par les automorphismes du prégraphe, ce qui
reviendrait à changer des relations entre $\varphi_0, \varphi$ via
$\lambda, \lambda^{\vee}$, en écrivant par ces relations
\[
\begin{cases}
\varphi = \check{\lambda} \varphi_0 = \varphi_0 \lambda^{-1} \quad [= (\lambda \varphi_0^{-1})^{-1}] \\
\text{donc } \text{\struck{$\varphi_0(\lambda)$}}
\begin{cases}
\check{\lambda} = \varphi_0(\lambda)^{-1} \\
\lambda = (\varphi_0^{-1}(\lambda^{\vee}))^{-1}
\end{cases} \\
\varphi_0(\lambda^{-1} s) = \check{\lambda} \varphi_0(s) \quad \text{ou} \quad \varphi_0(\lambda s) = \check{\lambda}^{-1} \varphi_0(s)
\end{cases}
\]
et qui donne pour \uncertain{contrainte}
\[
(\lambda, \check{\lambda})(\mathcal{R}) = \mathcal{R} \quad \text{i.e.} \quad
\forall (s, \check{s}) \in S \times S^{\vee},\ (s, \check{s}) \in \mathcal{R} \Longleftrightarrow (\lambda s, \check{\lambda} \check{s}) \in \mathcal{R} .
\]
\marginal{cette définition des $\lambda, \lambda^{\vee}$ en termes de
$\varphi_0, \varphi$ a aussi l'avantage de \uncertain{rendre} \ill{}
l'\uncertain{involution} symétrique $\lambda, \lambda^{\vee}$ qui
\uncertain{identifie} \ill{}}
Dans le cas de l'exemple linéaire, cela signifie que si on fait opérer les
scalaires\add{$\in \Lambda^{*}$} sur $M$ de la façon habituelle, — les
faire opérer sur $M^{\vee}$ par contragrédience (la « multiplication »
inhabituelle), ce qui a l'avantage

\page{210}
\note{Pagination de l'auteur : p. 108.}
de respecter la présence scalaire dans la relation d'incidence $\mathcal{R}$,
mais l'inconvénient que les applications ne commutent plus à l'op.\ de
$\Lambda^{*}$, sur $S$ et $\check{S}$, mais qu'il \struck{faut} écrire
$\varphi_0(\lambda s) = \check{\lambda}^{-1} \varphi_0(s)$
(au lieu de $\lambda \varphi_0(s)$). L'avantage d'avoir une loi de groupe
bien en évidence sur les homothéties, dans cette notation, \uncertain{passe}
peut-être préférablement. Mais il faut bien noter qu'en général, \emph{on n'a
pas} d'opération du groupe $Z$ sur l'ensemble $\Pi$ des « jauges » du
prégraphe $\Gamma$. \struck{On a bien défini}
\note{Le cadre qui suit est barré de plusieurs traits obliques.}
\[
\text{\struck{application}}\quad
\boxed{\text{\struck{$Z_S \times \Pi \to \Pi,\quad (\lambda, \varphi_0) \mapsto \varphi_0 \lambda^{-1} = \check{\lambda} \varphi_0$}}}
\]
Pour que $(\lambda, \check{\lambda}) \in \mathfrak{S}_S \times \mathfrak{S}_{\check{S}}$
définisse — à partir d'une jauge $\varphi_0 \in \Pi$ — une deuxième jauge
$\varphi$ $(= \check{\lambda} \varphi_0 = \varphi_0 \lambda^{-1})$ et qu'il en
soit ainsi, en plus de la condition que $(\lambda, \check{\lambda})$ soit un
automorphisme du prégraphe $\Gamma$, la condition qui \emph{dépend du choix
de $\varphi_0$})
\[
(65)\qquad \check{\lambda} \varphi_0 = \varphi_0 \lambda^{-1} \quad \text{i.e.} \quad \varphi_0 = \check{\lambda} \varphi_0 \lambda .
\]
Si on décide d'oublier $\check{\lambda}$ (p.\ ex.)\add{et le remplacer
par $\check{\lambda} = (\varphi_0^{-1}(\lambda))^{-1}$}, i.e.\ retenir que
$\lambda$ (et définir $\varphi$ par $\varphi_0 \lambda^{-1}$), on
\uncertain{conjugue} : la \struck{\ill{}} \uncertain{première} condition sur
$\lambda$ qui exprime \struck{la symétrie} fait que
\[
(\lambda, \check{\lambda}) = (\lambda, \varphi_0^{-1}(\lambda^{-1}))
\]
\struck{satisfasse la condition de symétrie} (1° \uncertain{inverse}
$\mathcal{R}$ ; 2°) ne définit pas \uncertain{a priori} un \emph{sous-groupe}
de $\mathfrak{S}_S$ ! En fait, si on exprime

\page{211}
cette condition sous la forme (29'),
\[
\forall s, t \in S,\quad \{s, t\} \in A_0 \Longleftrightarrow \{\lambda^{-1} s, \lambda t\} \in A_0 ,
\]
on voit ceci
\[
(66)\qquad
\begin{array}{l}
\text{si } \lambda, \mu \in \mathfrak{S}_S \text{ satisfont séparément à la condition} \\
\text{(cas lin.\ i.e.\ ce sont des homothéties pour le graphe } (S, A_0)\text{),} \\
\text{alors } \lambda\mu \text{ y satisfait si } \lambda\mu = \mu\lambda \\
\text{(cas d'un graphe \emph{séparé}).}
\end{array}
\]
Pour continuer le travail, on supposera donc donné un sous-groupe
\emph{commutatif}
\[
\text{\struck{\ill{}}}\qquad Z \subset \mathfrak{S}_S
\]
formé d'homothéties de $(S, A_0)$, et on considère \struck{dans}
\struck{les éléments} les sous-groupes correspondants\add{$Z_{\Gamma}$}
\struck{de $\mathfrak{S}_S \times \mathfrak{S}_{S^{\vee}}$}, de
$\mathrm{Aut}(S, S^{\vee}, \mathcal{R})$, on trouve : \struck{que les}
\[
(67)\qquad
\begin{cases}
\text{les } \varphi_\alpha : S \to \check{S} \text{ telles qu'on ait} \\
\quad \check{\lambda} \varphi_\alpha \lambda = \varphi_\alpha \text{ pour tout } \lambda \in Z_{\Gamma} \subset \mathrm{Aut}_S(\Gamma), \\
\text{forment un } \text{\struck{torseur sous le groupe } $Z_{\Gamma}$}\ \text{ensemble } \Pi_Z, \text{ sur} \\
\text{lequel le groupe } Z_{\Gamma} \text{ y opère \emph{librement} par l'opération} \text{\struck{\ill{}}} \\
\quad T_{\lambda, \check{\lambda}}\, \varphi = \check{\lambda} \varphi = \varphi \lambda^{-1} .
\end{cases}
\]
\marginal{On a supposé $\Pi_Z \neq \emptyset$, i.e.\ que $\varphi_0 \in \Pi_Z$,
i.e.\ \ill{} \ill{}}
En effet, \struck{pour} si $(\lambda, \check{\lambda}) \in Z$, — doit vérifier
d'abord que
\[
\varphi \in \Pi_Z \Longrightarrow \underbrace{T_{\lambda, \check{\lambda}}\, \varphi}_{\check{\lambda}\varphi} \in Z \quad \text{i.e.} \quad
\forall (\mu, \check{\mu}) \in Z,\ \text{on a}
\]
\[
\check{\mu}(\check{\lambda} \varphi) \mu = \check{\lambda} \check{\mu} \varphi \mu = (\check{\lambda} \varphi) .
\]
\note{L'auteur écrit « $\in Z$ » là où le sens demande « $\in \Pi_Z$ ».}
OK, le fait que $Z_{\Gamma}$ opère librement est clair, \struck{\ill{}} car
pour $\varphi_0, \varphi$ fixés, il y a qu'un seul $(\lambda, \check{\lambda})$ …
\struck{Ensuite le fait que si $\varphi_0 \in$}
On a encore que

\page{212}
\note{Pagination de l'auteur : p. 109. Page très chargée : de longues notes
marginales, écrites en diagonale dans la marge gauche, sont transcrites à
part à la fin de la page, autant qu'elles se lisent.}
cette opération est transitive — soit donc $\varphi_0 \in \Pi_Z$, disons, et
$\varphi \in \Pi_Z$, on peut écrire
\[
\varphi = \check{g}_0 \varphi_0 = \varphi_0 g_0^{-1} \quad \text{où } (g_0, \check{g}_0) \in \mathrm{Aut}(\Gamma)
\]
\uncertain{est-il dans} $Z_{\Gamma}$ ? \uncertain{Ainsi} dites \uncertain{maintenant}
\add{$\mathrm{Aut}\,\varphi$}, on doit exprimer le $\varphi \in \Pi_Z$ i.e.
\[
\underbrace{\check{\lambda} \varphi \lambda}_{} = \varphi \quad \forall (\lambda, \check{\lambda}) \in Z \quad \text{i.e., écrivant } \varphi = \varphi_0 g_0^{-1}
\]
\[
\text{\struck{\ill{}}}\qquad
\check{\lambda} \varphi_0 g_0^{-1} \lambda = \varphi_0 g_0^{-1} \quad \text{i.e.}
\]
\[
\underbrace{\check{\lambda} \varphi_0}_{\varphi_0 \lambda^{-1}} g_0^{-1}(\lambda) = \varphi_0 \quad \text{i.e.}
\qquad \text{\struck{$\varphi_0 = \lambda_0 \varphi_0 \lambda^{-1}$}}
\]
\[
\lambda^{-1} g_0^{-1}(\lambda) = 1, \quad \text{i.e.} \quad \lambda = g_0^{-1}(\lambda) \quad \text{i.e.} \quad g_0(\lambda) = \lambda
\]
\note{Ici $g_0(\lambda)$ désigne le conjugué $g_0 \lambda g_0^{-1}$.}
i.e.\ $g_0$ \struck{normalise} centralise tout $\lambda \in Z$, \uncertain{provenant de}
\struck{\ill{}} $Z$. Donc on trouve
\[
(68)\qquad
\begin{cases}
\text{l'ens.\ } \Pi_Z \text{ des jauges de } \Gamma = (S, S^{\vee}, \mathcal{R}) \text{ compatibles avec le groupe} \\
\text{d'automorphismes } Z \subset \mathrm{Aut}(\Gamma) \text{, i.e.\ satisfaisant } \check{\lambda}\varphi\lambda = \varphi\ \forall (\lambda, \check{\lambda}) \in Z_{\Gamma}, \\
\text{est un \struck{torseur} sous le \emph{centralisateur} de } Z \text{ dans } \mathrm{Aut}(\Gamma), \text{ opérant par} \\
\quad T_g \varphi_0 = g_{S^{\vee}} \varphi_0 = \varphi_0 g_S^{-1}
\end{cases}
\]
Ouf — il commence \uncertain{à} fallait de le savoir, pour \uncertain{des}
conclusions de cette petite \uncertain{sorcie} ! \marginal{faux}
\note{L'auteur a entouré « torseur » et l'a relié par un trait au mot
« faux » écrit plus bas.}

\ill{} tendre qu'il \uncertain{montre} c'est que l'application
\[
g \longmapsto g_{S}^{\vee} \varphi_0 = \varphi_0 g_S^{-1}
\]
est une \emph{bijection} de l'ensemble $\mathcal{G}(\varphi_0)$ de $\mathcal{G}$
\add{NB on a $\mathcal{G}(\varphi_0) = \{g \in \mathrm{Aut}(\Gamma) \mid \varphi_0(g_S) \check{g} = 1\}$, i.e.\ $\check{g} = \varphi_0 g_S^{-1}$}
formé des $g$ satisfaisant $g_{S}^{\vee} \varphi_0 = \varphi_0 g_S^{-1}$ i.e.\ qui
sont des $A_{0}$-homothéties (et qui \emph{ne forment pas de groupe} !) sur l'ens.\
$\Pi_Z$ des $Z$-jauges ! Mais il est vrai que $Z$ opère librement sur $\Pi_Z$.

\marginal{\emph{Lemme} Soit $\varphi_0 \in \Pi_Z$ une $Z$-jauge, et
$g \in \mathrm{Aut}(\Gamma)$, telle que l'on ait $g_{S}^{\vee}\varphi_0 = \varphi_0 g_S^{-1}$
(i.e.\ $(g_S, g_{S^{\vee}})$ définit une $A_{\varphi_0}$-homothétie du
graphe \ill{}). Conditions \ill{} : a) $g_S^{\vee} \varphi = \varphi g_S^{-1}$
\ill{} pour $\mathcal{H}_\varphi$ i.e.\ \ill{} ; b) $g \in \mathrm{Cent}(Z)$
i.e.\ $g_S^{\vee}\varphi_0 = \varphi_0 g_S^{-1}$ \ill{} ; c)
$g \in$ \ill{}.}
\marginal{Il faut écrire la petite vérification que si $\varphi_0$,
$g_S^{\vee}\varphi_0 = \varphi_0 g_S^{-1}$ pour tous les
$\varphi \in \Pi$, \ill{} $g \in \mathcal{G}$).}
\marginal{Il est donc \uncertain{prudent} d'étendre un \ill{} \uncertain{pendant}
\ill{} l'ens.\ $\Pi_Z$, qui n'est pas bien compris et peut-être
\ill{} \ill{} pour \ill{} \ill{} ensemble $\pi$}

\page{213}
\marginal{Scholie}
En résumé, nous \struck{\ill{}} amenés à nous intéresser à la structure
suivante
\[
(69)\qquad
\begin{cases}
\text{a) Triple } \Gamma = (S, S^{\vee}, \mathcal{R}) \text{ de deux ens.\ } S, S^{\vee} \text{ et d'une relation } \mathcal{R} \subset S \times S^{\vee} \\
\text{b) Donnée d'un sous-groupe \emph{commutatif}} \\
\qquad Z \subset \mathrm{Aut}(\Gamma), \\
\quad \text{dont le centralisateur dans } \mathrm{Aut}(\Gamma) \text{ est noté } \mathcal{G}.
\end{cases}
\]
\marginal{dans le cas de l'exemple linéaire, avec $M^{*}$, $M^{\vee *}$, … ,
$\mathcal{G}$ contient le groupe $GL_{\Lambda}(M)$}
\note{Les deux lectures de l'exemple linéaire dans la marge gauche sont
incertaines.}
On s'intéresse aux \emph{$Z$-jauges} d'une telle structure, les applications
bijections
\[
(70)\qquad
\begin{cases}
\varphi : S \xrightarrow{\ \sim\ } S^{\vee} \text{ satisfaisant} \\
\text{1) } \text{\struck{\ill{}}}\ (\mathrm{id} \times \varphi)^{-1}(\mathcal{R}) \text{ est une relation \emph{symétrique} sur } S \\
\qquad \text{i.e.\ } \forall s, t \in S,\ (s, \varphi(t)) \in \mathcal{R} \Longleftrightarrow (t, \varphi(s)) \in \mathcal{R} \\
\text{2) } \forall (\lambda, \check{\lambda}) \in Z, \text{ on a} \\
\qquad \check{\lambda} \varphi \lambda = \varphi \quad \text{i.e.} \quad \check{\lambda} \varphi = \varphi \lambda^{-1}
\end{cases}
\]
Si $\mathcal{R}$ est bisignante, ce qui est le cas si on fait le donné
$(S, S^{\vee} \subset \mathfrak{P}(S))$, et $\mathcal{R}$ s'identifie : celles
des $\varphi$ satisfaisant (1) s'identifient aux structures de graphe
\add{strict} \struck{\ill{}} séparé, sur $S$ (sans exclure des boucles)
$A_0 \subset \mathfrak{P}_2(S)$, qui admettent $S^{\vee}$ comme ensemble des
étoiles, et pour lesquels les $\lambda \in Z_S$ (\uncertain{provenant} des
$\lambda \in Z$) sont des \emph{homothéties}. Les
\add{(En tous cas, il y a une action de $Z$ sur les)} $\varphi$ en question,
forment un \emph{torseur} \uncertain{non} sous $Z$, \uncertain{mais}) sous
$\mathcal{G} = \mathrm{Cent}_{\mathrm{Aut}(\Gamma)}(Z)$. Les $Z$-multi-graphes
\uncertain{envisagés} précédemment sont donc \uncertain{les}
sous-$Z$-torseurs de celui-ci, i.e.\ \struck{\ill{}} s'identifient aux points
de $\mathcal{G}/Z$.
\marginal{caduc ; cf.\ \uncertain{annotation} p.\ précédente. Mais $Z$ opère
librement sur $\Pi_Z$}

\page{214}
\note{Pagination de l'auteur : p. 110.}
Dans ce contexte général, si on part d'une « vraie » structure de graphe
(sans boucles, i.e.\ $A_0 \subset \mathfrak{P}_2(S)$) séparé, il n'est pas
clair qu'en la transformant par un élément de $Z$, faisant
\[
\text{donc } A_{\lambda} = \bigl\{\{s, \lambda t\} \in \mathfrak{P}(S) \bigm| s, t \in S \text{ t.q.\ } \{s, t\} \in A_0\bigr\},
\]
on trouve encore des vrais graphes — car il pourrait arriver qu'on ait
$t = \lambda s$ pour un $(s, t, \lambda) \in S \times S \times Z$, avec
$\{s, t\} \in A_0$.

Prenons p.\ ex.\ en $S^{\vee} = S$, $\mathcal{R} = $ diagonale, donc les jauges
(satisfaisant (70)(1)) sont les $\varphi \in \mathfrak{S}_S$ telles que la
relation \struck{\ill{}} $\varphi(t) = s$ en $s, t$ soit symétrique, i.e.\
équivalente à $\varphi(s) = t$ i.e.\ à $s = \varphi^{-1}(t)$, ce qui signifie
aussi
\[
(71)\qquad \varphi^2 = \mathrm{id} \qquad \text{i.e.\ } \varphi \text{ est une involution.}
\]
Les automorphismes $(\lambda, \check{\lambda})$ de $\Gamma = (S, \check{S}, \mathcal{R}_{\mathrm{diag}})$
sont les couples $(\lambda, \lambda)$ avec $\lambda \in \mathfrak{S}_S$, et
les conditions (70) cas \uncertain{s'écrivent}
\[
(72)\qquad \lambda \varphi \lambda = \varphi \quad \text{soit} \quad \varphi(\lambda) = \lambda^{-1}
\]
i.e.\ l'involution « renverse » $\lambda$. Si donc on se donne un sous-groupe
commutatif $Z$ de $\mathfrak{S}_S$, alors $\Pi_Z$ est formé des involutions
$\varphi \in \mathfrak{S}_S$ qui \struck{normalisent} $Z$, en y induisant
l'involution $\lambda \mapsto \lambda^{-1}$. Si \struck{\ill{}} p.\ ex.\
$Z$ est \uncertain{un} centralisateur dans $\mathfrak{S}_S$) \add{l'un sous-groupe
d'un torseur $L$} \struck{les $\varphi$ en question s'identifient aux
involutions de} $Z$, en \uncertain{donnant} \struck{les isomorphes} \ill{}
\struck{\ill{} involutions} \add{alors divisions à un} $L \simeq Z$, en
notant les involutions de la forme
\[
\varphi_{\lambda}(z) = \lambda - z, \qquad \lambda \in Z,\ z \in Z,
\]
\note{Ce passage, raturé et surchargé, n'est lu que partiellement.}
\struck{Ici la relation $s = \varphi_a(s)$ i.e.\ $z = a - z$}

\page{215}
on a
\[
(73)\qquad (z, \varphi_\lambda(z)) \in \mathcal{R} \Longleftrightarrow \varphi_\lambda(z) = z \quad \text{i.e.} \quad 2z = \lambda
\]
et \struck{rien n'empêche qu'il y ait} il y a des solutions\add{en $z$} de
cette équation pour certaines valeurs des paramètres (\struck{\ill{}} pour
$\lambda \in 2Z$ — NB on écrit additivement ici la loi de groupe), alors qu'il
est possible qu'il n'y en ait pas pour d'autres (dès lors que $2Z \neq Z$) —
i.e.\ qu'on peut \add{\uncertain{maintenant}} choisir une « origine »
$\varphi_{\lambda_0}$ telle qu'elle corresponde à une « vraie » structure de
graphe $A_{\lambda_0}$. L'exemple le plus économique serait
\struck{\ill{}} $Z = \mathbb{Z}/2\mathbb{Z}$,
…

Soit $\Pi$ un sous-$Z$-torseur du $\mathcal{G}$-torseur des $Z$-jauges du
prégraphe $\Gamma = (S, S^{\vee}, \Gamma)$.\note{Ainsi sur la page :
« $\Gamma = (S, S^{\vee}, \Gamma)$ » ; le troisième terme est sans doute
$\mathcal{R}$.} Ainsi $\Pi$ est formé de bijections
\[
\varphi : S \xrightarrow{\ \sim\ } \check{S}
\]
satisfaisant
\[
(74)\qquad \lambda_{\check{S}}\, \varphi(s) = \text{\struck{$\lambda_S$}}\ \varphi(\lambda_S^{-1}(s)), \qquad
\lambda_{S^{\vee}} \varphi = \varphi \circ \lambda_S^{-1}
\]
i.e.\ compatible avec les opérations de $Z$ sur $S$ et $\check{S}$, via
l'hom.
\[
(75)\qquad
\begin{array}{ccc}
\lambda & \longmapsto & \lambda^{-1} \\
Z & \longrightarrow & Z
\end{array}
\]
des groupes d'opérateurs. \struck{On peut donc considérer, si $\lambda \in Z$,
$\check{S}$ est \uncertain{déduit} de $S$ \uncertain{par} le $Z$-torseur
$\Pi$, pour l'opération de $Z$ sur $S$ : $T_\lambda \varphi = \lambda_{\check{S}} \varphi = \varphi \lambda_S^{-1}$}
\note{Les dernières lignes sont barrées de traits obliques et horizontaux ;
leur lecture est très incertaine.}

\page{216}
\note{Pagination de l'auteur : p. 111.}
Considérons l'application
\[
\begin{array}{rccl}
F : & \Pi \times S & \longrightarrow & \check{S} \\
& (\varphi, s) & \longmapsto & \varphi(s)
\end{array}
\]
on aura donc
\[
\underbrace{F(T_\lambda \varphi, s)}_{(\varphi \circ \lambda^{-1})(s) = \varphi(\lambda^{-1} s)} = F(\varphi, \lambda_S^{-1} s)
\]
ce qui montre que $F$ se factorise en une application canonique, qui est une
bijection d'ens.
\[
(76)\qquad \Pi \wedge_Z (S!) \xrightarrow{\ \sim\ } \check{S}, \qquad \varphi \wedge s \longmapsto \varphi(s)
\]
où $S!$ désigne le $Z$-ensemble déduit de $S$ en y remplaçant
l'\struck{\ill{}} opération donnée de $Z$ par l'opération « inverse »,
\add{« déduite »} qui s'en déduit par $\lambda \mapsto \lambda^{-1}$.
\struck{Dans $\Pi^{\vee}$ \ill{} \ill{} morphismes, \ill{} $\varphi$ \ill{}
\ill{} Comment les opérations de $Z$ sur $S$ et $\check{S}$}
\[
\underbrace{\lambda . (\varphi \wedge s)}_{\substack{(T_\lambda \varphi) \wedge s \,=\, \varphi \wedge \lambda_S^{-1} s \\ \text{par définition}}} \longrightarrow \lambda_{\check{S}}(\varphi(s))
\]
On voit que (76) est un isomorphisme de $Z$-ensembles. \struck{\ill{}} Est-ce
même un isomorphisme de $\mathcal{G}$-\struck{torseurs} ensembles ? (En notant
que $\mathcal{G}$ agit \struck{sur $S!$} par $Z$-automorphismes, il agit aussi
sur l'ensemble tordu $\Pi \wedge_Z (S!)$ par transport de structure,
\[
(77)\qquad \text{par } \quad g(\varphi \wedge s) = \varphi \wedge g_S(s)\ )
\]
On aura
\[
g(\varphi \wedge s) = \varphi \wedge g_S(s) \longmapsto \varphi(g_S(s)) \overset{?}{=} g_{\check{S}}\, \varphi(s) .
\]

\page{217}
\struck{La relation signifie donc que pour}
Ainsi
\[
(78)\qquad \Pi \wedge_Z (S!) \xrightarrow{\ \sim\ } S^{\vee} \text{ commute à un } g \in \mathcal{G} \text{ donné}
\]
si et seulement si pour tout $\varphi \in \Pi$ — on simplement pour \emph{un}
$\varphi_0 \in \Pi$ — on a $\varphi(g_S(s)) = g_{\check{S}}\, \varphi(s)$, i.e.\
$\varphi$ commute à $g$.

Donc, pour un sous-groupe donné $G \subset \mathcal{G} \subset \mathrm{Aut}(\Gamma)$
(on suppose aussi $Z \subset G$), l'isom.\ (76) est un isom.\ de $G$-ens.\
si pour \struck{tout} $\varphi \in \Pi$ — ou pour un $\varphi_0 \in \Pi$, c'est
kif-kif — $\varphi$ commute à l'action de $\mathcal{G}$. \struck{Le plus} Il y a
un plus grand tel groupe\add{$GL_{\Lambda}(M)$} $G$, qui \struck{sont} est
\add{donné} \struck{formé des}
\[
\text{\struck{$g \in \mathcal{G}$ qui}}
\]
\[
(79)\qquad \mathcal{G}_0 = \Bigl\{ g \in \mathrm{Aut}(S, S^{\vee}, \mathcal{R}) \Bigm|
\begin{array}{l}
\text{a) } g \text{ centralise } Z \text{ i.e.\ } g \in \mathcal{G} \\
\text{b) } \text{\struck{$\forall \varphi$}}\ \varphi \in \Pi,\ \varphi : S \simeq S^{\vee} \text{ commute} \\
\quad \text{à } g \text{ i.e.\ } \varphi(g_S s) = g_{\check{S}}(\varphi(s))
\end{array}
\Bigr\} .
\]
\marginal{NB Seconde des conditions entraîne la première : $g$ i.e.\ cette
condition de commutation au sens strict, avec la condition
$\check{g}\varphi = \varphi g^{-1}$ (au lieu de $\varphi g$) !}
Prenons $Z = 1$ dans $\mathcal{G} = \mathrm{Aut}(\Gamma)$. Alors $\Pi$ se
réduit à un élément, que on n'importe quelle jauge $\varphi_0$ du
$\Gamma = (S, S^{\vee}, \mathcal{R})$. Bien entendu, il n'est pas vrai, i.e.\
que $\varphi_0$ commute à l'action de $\mathcal{G}$ — p.\ ex.\ dans le cas de la
p.~110 $S = S^{\vee}$, $\mathcal{R} = $ diagonale, $\varphi_0$ est donc une
involution de $S$, tandis que $\mathcal{G} = \mathfrak{S}_S$ ; ici
$\mathcal{G}_0$ s'identifie au commutant de l'involution $\varphi_0$ dans
$\mathfrak{S}_E$.\note{« $\mathfrak{S}_E$ » sur la page ; sans doute pour
$\mathfrak{S}_S$.} \struck{Donc on a}

Il faut encore voir comment se transforment
\add{(\struck{\ill{}} signés)} \uncertain{une} structure de multigraphes,
à homothéties, rendues par… dans l'optique présente

\marginal{Exemple : \uncertain{dans} un \ill{} \ill{} idéal, pour \ill{}
\ill{} ce sont les $M^{*}$, $M^{\vee *}$ \ill{} $\mathcal{G}_0$ \ill{} \ill{}
le groupe $GL_{\Lambda}(M)$. Alors \ill{}}
\marginal{Exemple : $S$, $S^{\vee}$ \ill{} \ill{} \uncertain{(cf.\ \ill{})}
$\mathcal{G}_0$ est formé des $u \in GL_{\Lambda}(M)$ satisfaisant
$\check{u}\varphi_0 = \varphi_0 u$ i.e.\ $\check{u}\varphi_0 u^{-1} = \varphi_0$
\ill{}, i.e.\ $\det u = \pm 1$, i.e.\ \ill{} $\varphi_0 \det(u)$ \ill{}, donc
(79 bis) $\mathcal{G}_0 = GL^{\pm 1}_{\Lambda}(M)$ \ill{} \ill{} \ill{}}

\page{218}
\note{Pagination de l'auteur : p. 112.}
comme une structure
\[
(81)\qquad \Bigl(S, S^{\vee}, \mathcal{R} \subset S \times S^{\vee},\ Z \subset \mathrm{Aut}(\underbrace{S, \check{S}, \mathcal{R}}_{\Gamma}),\ \Pi \subset \mathrm{Bij}(S, S^{\vee})\Bigr)
\]
où, pour mémoire
\[
(82)\qquad
\begin{cases}
\text{a) } S, S^{\vee} \text{ deux ensembles, } \mathcal{R} \subset S \times S^{\vee} \text{ une relation} \\
\text{b) } Z \text{ un sous-groupe \emph{commutatif} de } \mathrm{Aut}(\Gamma) \\
\text{c) } \Pi \text{ un ensemble de bijections } S \to S^{\vee}, \text{ satisfaisant} \\
\quad \text{(1) condition de symétrie (cf.\ (70)),} \\
\quad \text{(2) } \lambda_{S^{\vee}} \varphi = \varphi \lambda_S^{-1} \text{ i.e.\ } \lambda_{S^{\vee}} \varphi \lambda_S = \varphi \quad \forall \lambda \in Z, \\
\quad \text{(3) } \Pi \text{ est un sous-}Z\text{-torseur de } \Pi_Z
\end{cases}
\]
\struck{(cette condition} [NB les $\varphi$ satisfaisant (1), (2) forment un
\struck{torseur sous} \add{$\mathcal{G}$-}\uncertain{torseur} sous
$\mathcal{G} = \mathrm{Cent}_{\mathrm{Aut}(\Gamma)}(Z)$, pour
$T_g \varphi = g_{S^{\vee}} \varphi = \varphi g_S^{-1}$]
\struck{\ill{} \ill{}}
\note{Les crochets sur ce point sont entourés d'additions interlinéaires
partiellement illisibles.}
\marginal{la condition (1), sous les données a) b) c), signifie que
$\mathcal{R}$ soit séparée (bisignante), \ill{} \ill{} \ill{} les structures
de graphes \ill{} strict séparés (sans boucles) ; la condition (2) que le
prégraphe \ill{} un groupe d'homothéties pour ces structures de graphes}
\marginal{NB Il suffit de poser (2) pour un $\varphi_0 \in \Pi$, pour \ill{}}

Si on considère l'objet
\[
(83)\qquad \bigl(S' = S^{\vee},\ S'^{\vee} = S,\ \mathcal{R}' = {}^{s}\mathcal{R},\ Z' = {}^{s}Z,\ \Pi' = \Pi^{-1} = \{\varphi^{-1} \mid \varphi \in \Pi\}\bigr),
\]
on voit qu'elles satisfont aux conditions (82) que le système \uncertain{déjà
présenté} (80).\note{Le renvoi « (80) » ne correspond à aucun numéro lu dans
ce lot ; l'auteur passe de (79 bis) à (81).} Notons que la bijection
canonique
\[
(84)\qquad \varphi \longmapsto \varphi^{-1} \qquad \Pi \xrightarrow{\ \sim\ } \Pi'
\]
satisfait
\[
(85)\qquad \underbrace{(T_\lambda \varphi)^{-1}}_{\lambda_{S^{\vee}} \varphi} = \underbrace{T'_{\lambda'}(\varphi^{-1})}_{\varphi^{-1} \lambda_{S'}^{-1}}
\qquad \text{pour l'iso } \lambda \mapsto \lambda',\ Z \to Z' \text{ \uncertain{précisé},}
\]
\marginal{dans cette formule, l'identification $\lambda' = (\lambda_{S^{\vee}}, \lambda_S)$ si $\lambda = (\lambda_S, \lambda_{S^{\vee}})$}
i.e.\ (84) est \struck{\ill{}} un iso de $Z$-ensembles ! Si on considère
l'iso déduit de (76) en \struck{\ill{}} faisant le twist par
\struck{$\Pi$} \add{$\Lambda(\varphi)$} le torseur inverse de
\struck{$\Pi$}
\[
S! \simeq \Pi^{\mathrm{inv}} \wedge_Z \check{S}
\]
ou ce qui revient au même, en passant\note{lecture incertaine} ; l'action
opposée de $Z$ sur les deux membres
\[
(86)\qquad S \simeq \Pi \wedge_Z (\check{S}!) ,
\]

\page{219}
qui résulte de ce que l'iso.\ (76),\add{des types} pour les systèmes
duals des systèmes envisagés…

\struck{\ill{} \ill{}}

\emph{Transport de structure.} Soit $g$ un automorph.\ de $(S, S^{\vee}, \mathcal{R})$,
et soit $\varphi : S \to S^{\vee}$ une jauge (il n'est pas question ici de $Z$,
pour l'instant). Alors $\varphi$ définit une structure de graphe $A_\varphi$, et
on a par transport de structure
\[
(87)\qquad
\begin{cases}
g(\varphi) \overset{\text{déf}}{=} g_{\check{S}} \circ \varphi \circ g_S^{-1} \\
g_S(A_\varphi) = A_{g(\varphi)}
\end{cases}
\qquad \text{dirai \ill{}.}
\]
\struck{$\varphi$ est une $g$-jauge, p.\ ex.\ pour \ill{}}
\struck{Quand on dispose d'un $Z$ ss-groupe commutatif de
$\mathrm{Aut}(\Gamma = (S, S^{\vee}, \mathcal{R}))$, et $g \in Z$, $\varphi_0 \in \Pi_Z$,
alors \ill{} $\mathcal{G}$, \ill{} \ill{} $\mathrm{Aut}(\Gamma)$), alors si
$g \in \mathcal{G}$ et \ill{} $Z$-jauge $\varphi$ — i.e.\ satisfaisant à
$\lambda_{\check{S}} \varphi = \varphi \lambda_S^{-1}$ $\forall \lambda \in Z$,
\ill{} \ill{} \ill{} $\mathcal{G}$-torseur \ill{}}
\[
(88)\qquad T_g \varphi_0 = g_{\check{S}} \varphi_0 = \varphi_0 g_S^{-1}
\qquad \text{\struck{et $g \in \mathrm{Aut}(\Gamma)$}}
\]
\struck{(NB pour \ill{} $\varphi$ \ill{} $Z$ \ill{} \ill{} i.e.\ \ill{}}
\[
\text{\struck{$(89)\quad g_{\check{S}} \varphi = \varphi g_S^{-1} \Longleftrightarrow g \in \mathrm{Cent}_{\mathrm{Aut}(\Gamma)}(Z)\ (= \mathcal{G})$}}
\]
\note{Les lignes (88)–(89) sont barrées de hachures obliques ; la lecture des
parties biffées est très incertaine.}
Donc, on a alors, pour $g \in \mathcal{G}$
\[
(90)\qquad \forall g \in Z,\ \varphi_0 \in \underbrace{\Pi_Z}_{\text{ens.\ des } Z\text{-jauges}}
\qquad \underbrace{g(\varphi)}_{\text{cf.\ (87)}} = T_{g^2}\varphi .
\]
Plus généralement (pour mémoire) si $\Pi$ est un sous-$Z$-tors.\ de $\Pi_Z$,
\[
(91)\qquad \text{si } g \in \mathcal{G} = \mathrm{Cent}_{\mathrm{Aut}(\Gamma)}(Z), \text{ alors conditions équivalentes}
\]
\[
\begin{cases}
\text{a) } g_{\check{S}} \varphi = \varphi g_S^{-1} \text{ pour \emph{un} } \varphi \in \Pi \ (\text{disons, } \varphi_0) \\
\text{b) } g_{\check{S}} \varphi = \varphi g_S^{-1} \text{ pour tout } \varphi \in \Pi
\end{cases}
\]
(NB relation qui s'écrit $\varphi(g_S) = (g_{S^{\vee}})^{-1}$ !) et alors on a
\[
(92)\qquad g(\varphi) = T_{g^2}\varphi \qquad \text{pour toute } \varphi \in \Pi .
\]
\marginal{NB Correspondent \uncertain{aux} homothéties pour les structures
$A_\varphi$ ($\varphi \in \Pi$). Les $g$ formant \uncertain{un} ensemble
\uncertain{stable} \ill{} (\uncertain{indépendant} de \ill{}), \ill{} dans
\ill{} sous $Z$}

\page{220}
\note{Pagination de l'auteur : p. 113.}
\struck{(4)} \emph{Retour sur la condition d'isomorphie.} Soit
\[
(93)\qquad \mathcal{A} = \mathrm{Aut}(\Gamma), \qquad \Gamma = (S, S^{\vee}, \mathcal{R})
\]
\marginal{(on suppose $\mathcal{R}$ séparée en $\check{S}$)}
\[
(94)\qquad Z \subset \mathcal{A} \ \text{un sous-groupe commutatif donné,}
\]
d'où \struck{$Z$} $\Pi_Z \subset \mathrm{Bij}(S, \check{S})$, sur lequel $Z$
opère librement (a priori $\Pi_Z$ peut être vide, mais si),
\[
(95)\qquad \text{une sous-}Z\text{-torseur } \Pi \text{ de } \Pi_Z .
\]
On pose
\[
(96)\qquad \mathcal{G} = \mathrm{Cent}_{\mathcal{A}}(Z) .
\]
\struck{$(98)\quad \mathcal{G}' = \{g \in \mathcal{G}$ centralisant $Z\}$}

\emph{Lemme.} Soit $f \in \mathfrak{S}_S$. (Conditions équivalentes :
(i) $\exists\, \varphi_0 \in \Pi$ tel que $f(A_{\varphi_0})$ de la forme
$A_\varphi$ avec $\varphi \in \Pi$ ;

(i bis) pour \emph{tout} $\varphi_0 \in \Pi$, on a $f(A_{\varphi_0})$ de la
forme $A_\varphi$ pour $\varphi = f(\varphi_0)$ convenable — i.e.\
$f(\Pi) = \Pi$ ;

\uncertain{chaque} fois en structure)

\struck{(ii) $\exists f \in$} \struck{L'ensemble des $f$ forment une sous}
Les $f$ satisfaisant ces conditions sont $\in \mathcal{G}$, et forment un
sous-groupe $G$ de $\mathcal{G}$, contenant les groupes
\[
\text{\struck{$\mathcal{A}$}}\ \mathcal{G}_0 = \mathrm{Aut}(S, A_{\varphi})\ \text{\struck{$\Gamma$}}
\]
(indépendant du choix de $\varphi$) et $Z$. \struck{(\ill{} $\varphi \in \Pi$, \ill{})}
(\uncertain{au} $Z$ \uncertain{près}) Le sous-groupe des automorphismes de la
structure de multigraphe : homothéties $(S, S^{\vee}, \mathcal{R}, Z, \Pi)$
jouissant l'identité sur $Z$. On a un diagramme (97), où la ligne est exacte
\note{La flèche verticale porte un crochet d'inclusion ; elle est transcrite
comme l'inclusion $Z \hookrightarrow G$, ce que le texte (« contenant … $Z$ »)
indique.}

\begin{tikzcd}
1 \arrow[r] & \mathcal{G}_0 \arrow[r] & G \arrow[r, "\delta"] & Z \\
 & & Z \arrow[u, hook] \arrow[ur, "\lambda \mapsto \lambda^2"'] &
\end{tikzcd}

\note{L'argument se poursuit au-delà de ce lot (diagramme (97) et sa
discussion).}

\end{document}
